% !TeX root = ../../Book2.tex
\chapter{\texorpdfstring{\(L^1\)}{L¹} Fourier 变换与卷积}
\label{b2:ch:03}
\BookChapterNumber{b2:ch:03}
% Axler2020Measure Section11C：本轮实读印刷363--367、374--377页。
% 本章改用d维Gauss正则化组织反演；先证变换公式，不预设反演或Plancherel。
\begin{chapterabstract}
本章把周期函数的离散频率推广到欧氏空间上的连续频率。可积性首先保证变换有定义，Riemann--Lebesgue 引理进一步给出连续性与无穷远衰减。卷积将空间中的平移平均转化为频率乘法，Young 不等式负责控制相应范数。通过逼近恒等与 Gauss 核，反演公式从合法的绝对收敛积分中逐步建立；最后区分普通反演、Lebesgue 点处的正则化反演和仅依赖零变换的唯一性。
\end{chapterabstract}
\begin{learninggoals}
\begin{itemize}
\item\label{b2:item:03-goal-normalization}按统一的 \(2\pi\) 约定计算平移、调制、线性变换及微分公式，说明每个公式的可积条件。
\item\label{b2:item:03-goal-decay}理解固定函数的频率衰减与整个函数族的统一衰减有何区别。
\item\label{b2:item:03-goal-convolution}证明卷积的存在、等价类良定义及指数估计，核对换序所需的绝对可积性。
\item\label{b2:item:03-goal-inversion}用 Gauss 正则化建立反演和唯一性，辨认函数代表、连续点及 Lebesgue 点的作用。
\end{itemize}
\end{learninggoals}

本章在 \(\mathbb R^d\) 上使用 Lebesgue 测度。第一卷的连续紧支集稠密性、有限指数的平移连续性及 Lebesgue 点理论，在实际调用处给出编号；多元换元和积分号下求导沿用\chapref{b2:ch:01}。卷积的基本估计在本章按所需的一般指数完整推导。

\ProseParagraphBreak

先读\secref{b2:sec:0301}的归一化和运算规则，再把\secref{b2:sec:0302}的衰减与\secref{b2:sec:0303}的乘法对应连起来。\secref{b2:sec:0304}所构造的核既是逼近工具，也是\secref{b2:sec:0305}反演证明中的收敛控制。Gauss 变换的证明仅用常微分方程与实积分，因此不会借用尚未建立的复积分公式。

\ProseParagraphBreak

Sheldon Axler 的《Measure, Integration \& Real Analysis》\cite[Section 11C]{Axler2020Measure}采用一维 Poisson 核讲解类似的反演路线，可与本章的 Gauss 核方法比较。Elias M. Stein 与 Rami Shakarchi 的《Fourier Analysis: An Introduction》\cite{Stein2003Fourier}可作为从级数过渡到变换的伴读；本章的空间维数、符号及反演条件以所写陈述为准。

% !TeX root = ../../Book2.tex
\section{Fourier 变换}
\label{b2:sec:0301}

周期函数的频率由整数标记。若函数定义在整个欧氏空间上而不再周期重复，频率也变为连续变量。离散系数的求和于是让位于积分，但可积性、归一化与极限交换仍须逐项检查。

\subsection{定义与规范}
\label{b2:subsec:030101}

\begin{definition}\label{b2:def:03-fourier-transform}
设 \(d\) 为正整数，\(f\in L^1(\mathbb R^d;\mathbb C)\)。函数
\begin{equation}\label{b2:eq:03-fourier-definition}
 \widehat f(\xi)\coloneqq\mathcal Ff(\xi)
 =\int_{\mathbb R^d}f(x)\cdot e^{-2\pi\mathrm{i}x\cdot\xi}\,\mathrm{d}x
\end{equation}
称为 \(f\) 的\term[fourier-bianhuan]{Fourier 变换}[Fourier transform]。对 \(g\in L^1(\mathbb R^d)\)，另记
\[
 \check g(x)\coloneqq\int_{\mathbb R^d}g(\xi)\cdot e^{2\pi\mathrm{i}x\cdot\xi}\,\mathrm{d}\xi.
\]
此时只定义了正号积分；它何时反演 \(\mathcal F\)，将在\secref{b2:sec:0305}证明。
\end{definition}
\symidx{\(\mathcal Ff=\widehat f\)：欧氏空间 Fourier 变换}
\symidx{\(\check g\)：正号 Fourier 积分}

积分对每个 \(\xi\) 绝对收敛，且 \(|\widehat f(\xi)|\leq\|f\|_1\)，所以变换不依赖 \(f\) 的代表。这里 \(x\cdot\xi\) 是欧氏内积，虚数单位写作 \(\mathrm{i}\)。采用这个 \(2\pi\) 约定后，后面的反演公式没有额外常数，平方范数也将保持不变。若使用 \(e^{-\mathrm{i}x\cdot\xi}\) 的其他教材，需要同时转换频率尺度，不能只删去指数中的 \(2\pi\)。

\ProseParagraphBreak

在一维，取 \(a>0\)，区间示性函数给出一个可直接计算的模型：
\[
 \widehat{\mathbf1_{[-a,a]}}(\xi)=\dfrac{\sin(2\pi a\xi)}{\pi\xi}\quad(\xi\ne0),
 \quad \widehat{\mathbf1_{[-a,a]}}(0)=2a.
\]
原函数紧支集，并不使其变换紧支集；跳跃还使变换仅按 \(1/|\xi|\) 的尺度衰减。相反，对 \(a>0\)，分成正、负半轴积分可得
\[
 \mathcal F(e^{-a|x|})(\xi)=\dfrac{2a}{a^2+4\pi^2\xi^2}.
\]
这些公式同时提供了后面检验常数和反演条件的例子。

上面的连续频率可以用有限区间的级数作具体比较。在一维，设 \(f\) 连续且支撑在 \([-A,A]\)，并取 \(L>2A\)。将 \(f\) 在 \([-L/2,L/2]\) 上周期延拓，长度 \(L\) 的归一化 Fourier 系数为
\[
 c_L(k)=\frac1L\int_{-L/2}^{L/2}f(x)\cdot e^{-2\pi\mathrm i kx/L}\,\mathrm dx
       =\frac1L\widehat f(k/L).
\]
因此保留 \(|k|\leq RL\) 的有限频率和是
\[
 \frac1L\sum_{|k|\leq\lfloor RL\rfloor}
 \widehat f(k/L)\cdot e^{2\pi\mathrm i xk/L}.
\]
对固定 \(R>0\) 与 \(x\)，变换定义及控制收敛给出被求和函数在 \([-R,R]\) 连续；这一个和在 \(L\to\infty\) 时是频率步长 \(1/L\) 的 Riemann 和，趋于 \(\int_{-R}^R\widehat f(\xi)e^{2\pi\mathrm i x\xi}\,\mathrm d\xi\)。有限频带的极限只解释了积分为何自然出现；随后再令 \(R\to\infty\) 需要单独的收敛依据，不能从这个计算直接推出反演。归一化中的 \(1/L\) 正好成为频率积分元，所以采用 \(2\pi\) 后不再出现额外常数。

\subsection{平移}
\label{b2:subsec:030102}

对 \(f\in L^1(\mathbb R^d;\mathbb C)\)、\(h\in\mathbb R^d\)，沿用第一卷的平移约定 \(\tau_hf(x)\coloneqq f(x-h)\)。变量替换 \(y=x-h\) 给出
\[
 \widehat{\tau_hf}(\xi)=e^{-2\pi\mathrm{i}h\cdot\xi}\cdot\widehat f(\xi).
\]
平移保持 \(L^1\) 范数，故两侧都有定义。空间中的位置变化只改变频率中的相位，不改变 \(|\widehat f|\)。因而仅保留变换的模，通常不能恢复函数所在的位置。
\symidx{\(\tau_hf\)：平移，\(\tau_hf(x)=f(x-h)\)}

\subsection{调制}
\label{b2:subsec:030103}

对 \(\eta\in\mathbb R^d\)、\(f\in L^1(\mathbb R^d;\mathbb C)\)，算子
\[
 M_\eta f(x)\coloneqq e^{2\pi\mathrm{i}\eta\cdot x}\cdot f(x)
\]
称为\term[tiaozhi]{调制}[modulation]。合并指数便得
\[
 \widehat{M_\eta f}(\xi)=\widehat f(\xi-\eta).
\]
因此空间中的调制将频率整体平移。两种操作的次序要保留：\(\tau_hM_\eta=e^{-2\pi\mathrm{i}\eta\cdot h}M_\eta\tau_h\)，一般并不交换。
\symidx{\(M_\eta\)：调制算子}

\subsection{伸缩}
\label{b2:subsec:030104}

\begin{proposition}\label{b2:prop:03-linear-change}
设 \(d\) 为正整数，\(f\in L^1(\mathbb R^d;\mathbb C)\)，\(A\) 是可逆实 \(d\times d\) 矩阵，\(g(x)=f(Ax)\)，则
\[
 \widehat g(\xi)=|\det A|^{-1}\cdot\widehat f(A^{-\mathsf T}\xi).
\]
特别地，对 \(a>0\)，若 \(f_a(x)\coloneqq a^{-d}\cdot f(x/a)\)，则 \(\widehat{f_a}(\xi)=\widehat f(a\xi)\)。
\end{proposition}
\begin{proof}
在定义积分中令 \(y=Ax\)。体积元为 \(|\det A|^{-1}\mathrm{d}y\)，而 \(x\cdot\xi=y\cdot A^{-\mathsf T}\xi\)。绝对可积性由同一换元公式保证；取 \(A=a^{-1}I\) 再乘以归一化因子，即得后一式。
\end{proof}

一个在空间中压缩到宽度 \(a\) 的单位质量函数，其变换扩展到约 \(a^{-1}\) 的频率范围。这是尺度互换的初步形式；此处只有精确换元公式，还没有给出一般的测不准估计。


\begin{proposition}\label{b2:prop:03-symmetries}
设 \(d\) 为正整数，\(f\in L^1(\mathbb R^d;\mathbb C)\)，则 \(\widehat{\overline f}(\xi)=\overline{\widehat f(-\xi)}\)。若 \(f\) 实值且 \(f(-x)=f(x)\) 几乎处处，则其变换为实值偶函数。若对每个实正交矩阵 \(Q\) 都有 \(f(Qx)=f(x)\) 几乎处处，则 \(\widehat f(Q\xi)=\widehat f(\xi)\) 对每个 \(Q\) 及所有 \(\xi\in\mathbb R^d\) 成立。
\end{proposition}
\begin{proof}
第一式由对定义积分取共轭得到。偶性由反射换元得到，再与第一式组合即知实偶函数的变换为实值。对正交矩阵 \(Q\)，线性换元公式中的绝对行列式为 \(1\)、逆转置为 \(Q\)。若 \(f(Qx)=f(x)\)，便有 \(\widehat f(Q\xi)=\widehat f(\xi)\)。
\end{proof}

例如矩形 \(Q=\prod_j[-a_j,a_j]\) 的变换为各坐标区间变换的乘积，频率中的零点形成相应的超平面。球的示性函数则有旋转不变的变换；但旋转不变性并不提供其径向表达式，还须计算角积分。这里先由对称性确定变换的结构，具体特殊函数留待附录处理。

\ProseParagraphBreak

反射与共轭也是校验计算的工具。若一个实值偶函数被算出了非实的变换，或径向函数被算出了依赖坐标方向的变换，应先回查相位和换元。另一方面，实值函数不一定有实值变换：只有共轭对称是普遍结论，偶性是额外条件。

\subsection{微分与乘法}
\label{b2:subsec:030105}

\begin{proposition}\label{b2:prop:03-moments}
设 \(d\) 为正整数，\(m\) 为非负整数。若 \(f\in L^1(\mathbb R^d;\mathbb C)\) 且 \(|x|^m f\in L^1(\mathbb R^d)\)，则 \(\widehat f\in C^m(\mathbb R^d;\mathbb C)\)，并对 \(|\alpha|\leq m\) 有
\[
 \partial_\xi^\alpha\widehat f(\xi)
 =(-2\pi\mathrm{i})^{|\alpha|}\mathcal F(x^\alpha f)(\xi).
\]
单独考虑一个坐标 \(1\leq j\leq d\) 时，只要 \(f,x_jf\in L^1(\mathbb R^d;\mathbb C)\)，就有连续的偏导数 \(\partial_{\xi_j}\widehat f=-2\pi\mathrm{i}\mathcal F(x_jf)\)，无须假设其他坐标的一阶矩可积。
\end{proposition}
\begin{proof}
一阶差商中指数之差除以步长的模不超过 \(2\pi|x_j|\)。因此 \(|x_jf|\) 提供共同的可积控制，可以把差商极限移入积分。导数的连续性由相同控制下的控制收敛得到，这就证明了仅需一个坐标矩的结论。若具有 \(m\) 阶绝对矩，迭代时每次增加一个坐标因子；所有不超过 \(m\) 阶的单项式都由常数乘 \(1+|x|^m\) 控制，故每一步都合法。
\end{proof}

\begin{proposition}\label{b2:prop:03-derivative}
设 \(d\) 为正整数，\(1\leq j\leq d\)。若 \(f\in C^1(\mathbb R^d;\mathbb C)\)，且 \(f,\partial_jf\in L^1(\mathbb R^d)\)，则
\[
 \mathcal F(\partial_jf)(\xi)=2\pi\mathrm{i}\xi_j\widehat f(\xi).
\]
若 \(f\in L^1(\mathbb R^d;\mathbb C)\) 且其第 \(j\) 个弱导数属于 \(L^1(\mathbb R^d;\mathbb C)\)，公式也成立。
\end{proposition}
\begin{proof}
取光滑紧支集函数 \(\chi\)，在单位球上等于 \(1\)，并令 \(\chi_R(x)=\chi(x/R)\)。先对紧支集函数 \(\chi_Rf\) 分部积分，得到
\[
 \mathcal F(\chi_R\partial_jf)(\xi)
 +\mathcal F(f\partial_j\chi_R)(\xi)
 =2\pi\mathrm{i}\xi_j\mathcal F(\chi_Rf)(\xi).
\]
第一项及右侧用控制收敛取极限；中间误差的绝对值至多为 \(R^{-1}\|\partial_j\chi\|_\infty\cdot\|f\|_1\)，故趋零。弱导数情形直接把 \(\chi_Re^{-2\pi\mathrm{i}x\cdot\xi}\) 代入弱导数的测试恒等式，得到同一公式和估计。
\end{proof}

于是空间中的可积导数产生频率衰减，空间中的可积矩产生频率光滑性。两类假设各控制一个不同方向，不能把“变换连续”误认为原函数已经可微。

% !TeX root = ../../Book2.tex
\section{Riemann--Lebesgue 引理}
\label{b2:sec:0302}

周期情形中 Fourier 系数趋零；在欧氏空间中，对应的问题是频率沿任意方向走向无穷时，变换是否趋零。单纯观察某个积分公式不足以说明所有方向，必须用可积逼近将有限计算推广。

\subsection{基本衰减}
\label{b2:subsec:030201}

\begin{theorem}[Riemann--Lebesgue]\label{b2:thm:03-riemann-lebesgue}
Fourier 变换将 \(L^1(\mathbb R^d)\) 连续地映入 \(C_0(\mathbb R^d)\)，且
\[
 \|\widehat f\|_\infty\leq\|f\|_1.
\]
这里 \(C_0\) 表示在无穷远处趋零的连续函数空间。
\end{theorem}
\begin{proof}
对有界矩形 \(Q=\prod_{j=1}^d[a_j,b_j]\)，Fubini 把 \(\widehat{\mathbf1_Q}\) 分解成一维区间变换的乘积。每个因子连续且有界；当 \(|\xi|\to\infty\) 时，至少一个坐标的绝对值不小于 \(|\xi|/\sqrt d\)，相应因子趋零，所以整个乘积趋零。

有限个矩形示性函数的线性组合在 \(L^1\) 中稠密。具体可先用第一卷命题 10.2.4 的连续紧支集稠密性逼近 \(f\)，再把包含其支集的一个大矩形细分成网格；一致连续性使每个网格内用常数代替函数的总积分误差趋零。记所得阶梯函数为 \(q\)，则
\[
 \|\widehat f-\widehat q\|_\infty\leq\|f-q\|_1.
\]
因此 \(\widehat f\) 是 \(C_0\) 函数的一致极限，仍属于 \(C_0\)。范数界直接来自积分三角不等式。
\end{proof}

该结论沿用\subsecref{b2:subsec:020202}的 Riemann--Lebesgue 名称。对 \(f_n\to f\) 于 \(L^1\)，变换在全部频率上一致收敛，而不只是每个固定频率分别收敛。

\subsection{连续性}
\label{b2:subsec:030202}

实际上 \(\widehat f\) 一致连续。定义式给出
\[
 |\widehat f(\xi+h)-\widehat f(\xi)|
 \leq\int|f(x)|\cdot|e^{-2\pi\mathrm{i}x\cdot h}-1|\,\mathrm{d}x.
\]
右侧不依赖 \(\xi\)，并由控制收敛趋零。若又有 \(\int|x|\cdot|f(x)|\,\mathrm{d}x<\infty\)，则可将右侧进一步估计为 \(2\pi|h|\int|x|\cdot|f(x)|\,\mathrm{d}x\)，得到 Lipschitz 控制。

\ProseParagraphBreak

注意 \(L^1\) 假设并未强迫原函数有界或连续。例如 \(f(x)=|x|^{-a}\mathbf1_{\{|x|<1\}}\) 在 \(0<a<d\) 时可积，其变换仍连续且趋零。反过来，紧支集函数的变换也未必可积；区间示性函数的 \(1/|\xi|\) 尾部正好落在一维不可积的边界。

\subsection{无穷远行为}
\label{b2:subsec:030203}

\begin{proposition}\label{b2:prop:03-no-uniform-decay}
在 \(L^1(\mathbb R^d)\) 的单位球上，不存在只依赖 \(|\xi|\)、并在无穷远处趋零的共同上界 \(|\widehat f(\xi)|\leq r(|\xi|)\)。
\end{proposition}
\begin{proof}
取非负单位积分的 \(g\in L^1\)。对任意 \(\eta\)，函数 \(M_\eta g\) 的范数仍为 \(1\)，而 \(\widehat{M_\eta g}(\eta)=\widehat g(0)=1\)。令 \(|\eta|\) 任意大，便与共同趋零上界矛盾。
\end{proof}

Riemann--Lebesgue 引理逐个控制固定函数，并没有统一控制整个单位球。若增加导数范数或矩的上界，才可能提取定量信息。例如当所有 \(m\) 阶纯坐标弱导数都可积时，选取最大的频率坐标，由\cref{b2:prop:03-derivative}可得
\[
 |\widehat f(\xi)|\leq C_{d,m}|\xi|^{-m}\sum_{j=1}^d\|\partial_j^m f\|_1
 \quad(\xi\ne0).
\]
这也说明变换衰减的常数应明确由哪些原函数量控制。

% !TeX root = ../../Book2.tex
\section{卷积}
\label{b2:sec:0303}

频率上的乘法对应空间中的卷积。要使用这个对应关系，先须证明卷积本身存在，并明确它落在哪个函数空间。可积性在这里不仅是形式上的标签，也是交换积分的依据。

\subsection{卷积代数}
\label{b2:subsec:030301}

\begin{definition}\label{b2:def:03-convolution}
对可测函数 \(f,g\)，若积分绝对收敛，定义它们在 \(x\) 处的\term[juanji]{卷积}[convolution]为
\[
 (f*g)(x)\coloneqq\int_{\mathbb R^d}f(x-y)\cdot g(y)\,\mathrm{d}y.
\]
\end{definition}
\symidx{\(f*g\)：欧氏空间卷积}

\begin{proposition}\label{b2:prop:03-l1-algebra}
若 \(f,g\in L^1(\mathbb R^d)\)，则卷积几乎处处存在，属于 \(L^1\)，且
\[
 \|f*g\|_1\leq\|f\|_1\cdot\|g\|_1.
\]
卷积在 \(L^1\) 等价类上良定义，满足交换律、结合律和双线性。
\end{proposition}
\begin{proof}
可测函数可选取 Borel 代表；函数 \((x,y)\mapsto f(x-y)g(y)\) 因而联合可测。对绝对值用 Tonelli 和平移不变性，
\[
 \int\!\int|f(x-y)|\cdot|g(y)|\,\mathrm{d}y\,\mathrm{d}x
 =\|f\|_1\cdot\|g\|_1<\infty.
\]
这既保证几乎处处绝对收敛，又给出范数界。改变任一代表所产生的差，按同一估计为零，故等价类良定义。令 \(z=x-y\) 得交换律。三个函数的绝对值重积分等于三个范数之积，所以可用 Fubini 重排；两种卷积括号都等于
\[
 \int\!\int f(x-y-z)\cdot g(y)h(z)\,\mathrm{d}y\,\mathrm{d}z
\]
几乎处处。线性来自积分的线性。
\end{proof}

结合 \(L^1\) 的完备性，这使 \(L^1\) 成为卷积运算下的 Banach 代数。它没有单位元：若 \(e*f=f\) 对所有 \(f\in L^1\) 成立，后面将由变换证明 \(\widehat e\) 必须恒等于 \(1\)，与 Riemann--Lebesgue 引理矛盾；在分布空间中才会出现相应的 Dirac 单位。

\subsection{Young 不等式}
\label{b2:subsec:030302}

\begin{theorem}[Young，卷积形式]\label{b2:thm:03-young}
设 \(1\leq p,q,r\leq\infty\) 且 \(1/p+1/q=1+1/r\)，其中 \(1/\infty=0\)。若 \(f\in L^p\)、\(g\in L^q\)，则卷积几乎处处存在，并且
\[
 \|f*g\|_r\leq\|f\|_p\cdot\|g\|_q.
\]
\end{theorem}
\begin{proof}
若 \(r=\infty\)，指数 \(p,q\) 共轭，对每个 \(x\) 使用 Hölder 即得积分的绝对收敛与一致上界。

若 \(r<\infty\)，则 \(p,q\leq r\)。先对有界紧支集函数计算，把绝对值乘积拆成
\[
 |f(x-y)|\cdot|g(y)|
 =\bigl(|f(x-y)|^p\cdot|g(y)|^q\bigr)^{1/r}
    |f(x-y)|^{1-p/r}\cdot|g(y)|^{1-q/r}.
\]
Hölder 的三个倒指数为 \(1/r\)、\(1/p-1/r\)、\(1/q-1/r\)，和为 \(1\)。某个倒指数为零时相应因子为 \(1\)，将其省去。于是
\[
 |f*g(x)|^r
 \leq\left(\int |f(x-y)|^p\cdot|g(y)|^q\,\mathrm{d}y\right)
       \|f\|_p^{r-p}\cdot\|g\|_q^{r-q}.
\]
对 \(x\) 积分，Tonelli 给出右侧为 \(\|f\|_p^r\cdot\|g\|_q^r\)。一般函数先对绝对值作幅值和区域的递增截断，所得非负卷积由单调收敛及该统一范数界确定一个有限的 \(L^r\) 函数。这保证原复值积分几乎处处绝对收敛，积分三角不等式再给出所需估计。零范数的情形直接成立。
\end{proof}

这称为\term[young-juanji-budengshi]{Young 卷积不等式}[Young's convolution inequality]，与标量乘积的 Young 不等式区别在于对象和指数关系。常用情形是 \(L^1*L^p\subset L^p\)，以及 \(L^p*L^{p'}\subset L^\infty\)。后一种情形还可利用平移连续性推出连续代表。

\subsection{平移与卷积}
\label{b2:subsec:030303}

第一卷引理 13.4.5 已证明：若 \(1\leq p<\infty\)，则 \(\|\tau_hf-f\|_p\to0\)。该结论首先对连续紧支集函数成立，再由平移等距性和稠密性推广。无穷端点不在其中；区间示性函数的小平移在 \(L^\infty\) 中始终保持误差 \(1\)。

\begin{proposition}\label{b2:prop:03-convolution-continuity}
若 \(1<p<\infty\)、\(f\in L^p\)、\(g\in L^{p'}\)，则卷积具有有界一致连续的代表。\(L^1*L^\infty\) 的卷积也有这一性质。
\end{proposition}
\begin{proof}
Hölder 保证每个点的积分绝对收敛。由变量替换，
\[
 \|\tau_h(f*g)-f*g\|_\infty
 \leq\|\tau_hf-f\|_p\cdot\|g\|_{p'}\longrightarrow0.
\]
在 \(L^1*L^\infty\) 情形，将平移放在 \(L^1\) 因子上即可。
\end{proof}

连续性在这里由两个因子的整体平移行为取得，并非因为原函数本身连续。另一方面，\(L^1*L^1\) 一般只有可积性，不能仅凭卷积符号就称所得函数光滑。若某一因子为光滑紧支集函数，才可将导数移到该因子，并由可积控制证明结果光滑。


\begin{corollary}\label{b2:cor:03-convolution-c0}
若 \(1<p<\infty\)、\(f\in L^p\)、\(g\in L^{p'}\)，则 \(f*g\in C_0(\mathbb R^d)\)。
\end{corollary}
\begin{proof}
取连续紧支集函数 \(f_n\to f\) 于 \(L^p\)、\(g_n\to g\) 于 \(L^{p'}\)。可取两列范数有界。由 Hölder 型卷积界，
\[
 \|f*g-f_n*g_n\|_\infty
 \leq\|f-f_n\|_p\cdot\|g\|_{p'}
 +\|f_n\|_p\cdot\|g-g_n\|_{p'}\longrightarrow0.
\]
每个 \(f_n*g_n\) 连续且支撑在两个紧支集的和中，所以属于 \(C_0\)。一致极限仍属于 \(C_0\)。
\end{proof}

两个指数都有限，使两边都能用紧支集函数逼近。端点 \(L^1*L^\infty\) 不保证在无穷远消失：取 \(g=1\)，则卷积恒为 \(\int f\)。这与已经证明的有界一致连续性不矛盾。卷积的“平滑作用”应拆成可积性、连续性、无穷远行为和可微性分别判断，不能把这些性质合成一个没有条件的结论。

\subsection{Fourier 变换与卷积}
\label{b2:subsec:030304}

\begin{theorem}[卷积定理]\label{b2:thm:03-convolution-transform}
若 \(f,g\in L^1(\mathbb R^d)\)，则
\[
 \widehat{f*g}(\xi)=\widehat f(\xi)\cdot\widehat g(\xi).
\]
\end{theorem}
\begin{proof}
定义积分中的绝对值重积分等于 \(\|f\|_1\cdot\|g\|_1\)，故可交换次序。令 \(z=x-y\)，指数分解为 \(e^{-2\pi\mathrm{i}z\cdot\xi}e^{-2\pi\mathrm{i}y\cdot\xi}\)，两个积分分别为 \(\widehat f\) 与 \(\widehat g\)。
\end{proof}

这个\term[juanji-dingli]{卷积定理}[convolution theorem]将叠加平移的操作对角化。现在也可补足单位元不存在的论证：若 \(e*f=f\)，固定频率 \(\xi\) 后取单位质量非负函数的调制，使 \(\widehat f(\xi)=1\)，便有 \(\widehat e(\xi)=1\)。这与 \(\widehat e\in C_0\) 不相容。下一节改用一族核逐渐接近单位作用。

% !TeX root = ../../Book2.tex
\section{逼近恒等}
\label{b2:sec:0304}

圆周上，Fejér 核把质量集中到原点，同时保持积分为 \(1\)。欧氏空间中的同一思想还要控制无穷远尾部。只要这两种控制成立，核可以有符号，也不必紧支集。

\subsection{逼近恒等的条件}
\label{b2:subsec:030401}

\begin{definition}\label{b2:def:03-approximate-identity}
设 \((K_\varepsilon)_{\varepsilon>0}\subset L^1(\mathbb R^d)\)。若
\begin{enumerate}
\item\label{b2:item:03-identity-mass}\(\int K_\varepsilon=1\)；
\item\label{b2:item:03-identity-bound}\(\sup_{\varepsilon>0}\|K_\varepsilon\|_1<\infty\)；
\item\label{b2:item:03-identity-tail}对每个 \(\delta>0\)，\(\int_{|x|\geq\delta}|K_\varepsilon(x)|\,\mathrm{d}x\to0\) 当 \(\varepsilon\downarrow0\)，
\end{enumerate}
则称它为一族\term[bijin-hengdeng]{逼近恒等}[approximate identity]。第一卷定义 13.4.6 的固定核缩放是其中一种情形。
\end{definition}

确实，若 \(K\in L^1\)、\(\int K=1\)，取 \(K_\varepsilon(x)=\varepsilon^{-d}\cdot K(x/\varepsilon)\)，前三个量分别由换元化为 \(1\)、\(\|K\|_1\) 和 \(\int_{|z|\geq\delta/\varepsilon}|K(z)|\,\mathrm{d}z\)。最后一项因可积性趋零。固定核允许正负抵消；一般核族则必须另写统一 \(L^1\) 界，不能由积分等于 \(1\) 推出。

\subsection{\texorpdfstring{\(L^p\)}{Lᵖ} 逼近}
\label{b2:subsec:030402}

\begin{theorem}\label{b2:thm:03-approximate-identity}
设 \((K_\varepsilon)\) 满足上述条件。若 \(1\leq p<\infty\)、\(f\in L^p\)，则 \(\|K_\varepsilon*f-f\|_p\to0\)。若 \(f\) 有界且一致连续，则 \(K_\varepsilon*f\to f\) 一致成立。
\end{theorem}
\begin{proof}
记 \(A=\sup_\varepsilon\|K_\varepsilon\|_1\)。利用积分为 \(1\)，将误差写成 \(\int K_\varepsilon(y)(\tau_yf-f)\,\mathrm{d}y\)。对归一化正权 \(|K_\varepsilon|/\|K_\varepsilon\|_1\) 使用 Jensen，再对 \(x\) 积分，得到
\[
 \|K_\varepsilon*f-f\|_p^p
 \leq A^{p-1}\int|K_\varepsilon(y)|\cdot\|\tau_yf-f\|_p^p\,\mathrm{d}y.
\]
平移连续性使 \(|y|<\delta\) 时范数误差统一小于任给的正数；远处使用 \(\|\tau_yf-f\|_p\leq2\|f\|_p\)，再由核的尾部条件控制。先选 \(\delta\)，再令 \(\varepsilon\downarrow0\)，即可完成证明。有界一致连续情形直接对逐点误差作相同分割，近处用一致连续性，远处用 \(2\|f\|_\infty\)。
\end{proof}

有限 \(p\) 与一致连续性的假设不可随意省略。即使取非负紧支集核，跳跃函数也不能在 \(L^\infty\) 中被连续卷积一致逼近。另一方面，一族只满足总积分为 \(1\) 的核可以平移到无穷远，使其与固定函数的卷积远离原函数；质量的位置和绝对质量都要检查。

\ProseParagraphBreak

磨光函数是这种方法的常见实例。取非负 \(\rho\in C_c^\infty\)、\(\int\rho=1\)，则 \(\rho_\varepsilon*f\to f\) 于有限指数的 \(L^p\)，而卷积光滑。若还先对 \(f\) 做空间截断，就得到光滑紧支集逼近。对同时属于 \(L^1\) 和 \(L^2\) 的函数，可先用同一个有界紧支集截断，再把磨光尺度取得足够小，使两种范数误差同时趋零；下一章将使用这个共同稠密类。


\begin{proposition}\label{b2:prop:03-first-moment-error}
设 \(K\in L^1\)、\(\int K=1\)，且 \(M_1=\int|z|\cdot|K(z)|\,\mathrm{d}z<\infty\)。若 \(f\in C_c^1(\mathbb R^d)\)、\(1\leq p\leq\infty\)，则
\[
 \|K_\varepsilon*f-f\|_p
 \leq\varepsilon M_1\bigl\||\nabla f|\bigr\|_p.
\]
\end{proposition}
\begin{proof}
微积分基本定理给出
\[
 f(x-\varepsilon z)-f(x)
 =-\varepsilon\int_0^1 z\cdot\nabla f(x-s\varepsilon z)\,\mathrm{d}s.
\]
对绝对值以 \(\varepsilon|z||\nabla f|\) 控制，随后用权重 \(|z||K(z)|\,\mathrm{d}z\,\mathrm{d}s/M_1\) 的 Jensen 不等式并对 \(x\) 积分。平移保持梯度范数，得到误差的 \(p\) 次方至多为 \((\varepsilon M_1)^p\bigl\||\nabla f|\bigr\|_p^p\)。无穷端点直接在梯度上取上确界。\(M_1=0\) 对单位积分的普通可积核不可能发生；若只看零误差情形，也可单独处理。
\end{proof}

这个估计把定性逼近转成一次阶误差，并说明常数由核的一阶绝对矩和函数梯度控制。若核为偶函数，线性项的积分可能相消；但要把这种相消转成更高阶的范数误差，还需控制二阶余项。后面的 Gauss 核习题会具体计算这一改进。

\subsection{Gauss 核}
\label{b2:subsec:030403}

\begin{theorem}[Gauss 函数的 Fourier 变换]\label{b2:thm:03-gaussian}
令 \(G(x)=e^{-\pi|x|^2}\)。则 \(\int_{\mathbb R^d}G=1\)，且 \(\widehat G=G\)。因而
\[
 G_\varepsilon(x)\coloneqq\varepsilon^{-d}e^{-\pi|x|^2/\varepsilon^2},
 \quad \widehat {G_\varepsilon}(\xi)=e^{-\pi\varepsilon^2|\xi|^2}.
\]
\end{theorem}
\begin{proof}
先看一维。若 \(I=\int_{\mathbb R}e^{-\pi x^2}\,\mathrm{d}x\)，则由 Tonelli 和平面极坐标换元，
\[
 I^2=2\pi\int_0^\infty e^{-\pi r^2}r\,\mathrm{d}r=1,
\]
故 \(I=1\)。记 \(g=\widehat G\)。空间矩可积使 \(g\) 可微，而 \(G'=-2\pi xG\)；结合\cref{b2:prop:03-moments}和\cref{b2:prop:03-derivative}得到 \(g'(\xi)=-2\pi\xi g(\xi)\)。所以 \(e^{\pi\xi^2}g(\xi)\) 导数为零，由 \(g(0)=1\) 得 \(g(\xi)=e^{-\pi\xi^2}\)。

高维时 \(G(x)=\prod_j e^{-\pi x_j^2}\)，绝对可积性和 Fubini 将变换分解为一维结果的乘积。最后使用伸缩公式。
\end{proof}

函数 \(G_\varepsilon\) 称为归一化\term[gauss-he]{Gauss 核}[Gaussian kernel]。它有无限支集，但尾部因指数衰减而可控。其 Fourier 变换仍为同类函数，使空间磨光与频率截断之间出现一个可精确计算的对应。
\symidx{\(G_\varepsilon\)：单位质量 Gauss 核}

\subsection{Poisson 核}
\label{b2:subsec:030404}

在 \(\mathbb R^d\) 上另定义
\[
 P_t(x)\coloneqq c_d\dfrac{t}{(t^2+|x|^2)^{(d+1)/2}},
 \quad
 c_d^{-1}=\int_{\mathbb R^d}(1+|z|^2)^{-(d+1)/2}\,\mathrm{d}z.
\]
归一化常数有限且正：单位球内被积函数有界；在半径约 \(2^j\) 的环带上，积分至多为 \(C2^{-j}\)。这些核非负、积分为 \(1\)，并且 \(P_t(x)=t^{-d}P_1(x/t)\)，所以构成逼近恒等。它与上一章的圆周 Poisson 核定义域不同，本节用下标 \(t>0\) 区分欧氏空间的尺度。
\symidx{\(P_t\)：欧氏空间 Poisson 核}

\ProseParagraphBreak

一维时 \(c_1=1/\pi\)，所以 \(P_t(x)=t/(\pi(t^2+x^2))\)。由本章开头双边指数的直接计算，
\[
 \mathcal F(e^{-2\pi t|x|})(\xi)=P_t(\xi).
\]
此刻不能未经证明就反转变换；下一节的反演定理建立以后，才能推出 \(\widehat {P_t}(\xi)=e^{-2\pi t|\xi|}\)。Poisson 核的空间尾部仅按幂衰减，Gauss 核则按指数衰减；两者都逼近恒等，但对远处增长的容许范围不同。

% !TeX root = ../../Book2.tex
\section{Fourier 反演}
\label{b2:sec:0305}

反演的关键不是把指数的符号形式上改回来，而是说明两重振荡积分怎样得到原函数。直接交换两个没有绝对收敛保证的积分没有意义。我们先加一个可积的 Gauss 因子，使换序合法，再让这个因子逐渐趋于 \(1\)。

\subsection{光滑情形}
\label{b2:subsec:030501}

\begin{lemma}\label{b2:lem:03-gaussian-inversion}
若 \(f\in L^1(\mathbb R^d)\)、\(\varepsilon>0\)，则每个 \(x\) 满足
\begin{equation}\label{b2:eq:03-gaussian-inversion}
 \int_{\mathbb R^d}\widehat f(\xi)\cdot e^{-\pi\varepsilon^2|\xi|^2}e^{2\pi\mathrm{i}x\cdot\xi}\,\mathrm{d}\xi
 =(f*G_\varepsilon)(x).
\end{equation}
\end{lemma}
\begin{proof}
左侧绝对收敛，因为 \(\widehat f\) 有界而 Gauss 因子可积。代入变换定义，重积分绝对值的积分为 \(\|f\|_1\int e^{-\pi\varepsilon^2|\xi|^2}\,\mathrm{d}\xi<\infty\)，故可用 Fubini。内层正号积分由 Gauss 变换公式及偶性等于 \(G_\varepsilon(x-y)\)，从而得到右侧。
\end{proof}

\begin{proposition}\label{b2:prop:03-smooth-inversion}
若 \(f\in C_c^\infty(\mathbb R^d)\)，则 \(\widehat f\in L^1\)，并有处处成立的绝对收敛反演公式 \(f=\check {\widehat f}\)。
\end{proposition}
\begin{proof}
取整数 \(m>d\)。纯坐标高阶导数可积，故\cref{b2:prop:03-derivative}及最大坐标估计给出 \(|\widehat f(\xi)|\leq C(1+|\xi|)^{-m}\)，于是变换可积。在\eqref{b2:eq:03-gaussian-inversion}左侧以 \(|\widehat f|\) 控制，得到 \(\check {\widehat f}(x)\)；右侧由有界一致连续函数的逼近恒等结论一致趋于 \(f\)。
\end{proof}

此证明只用到已有的 Gauss 积分、导数变换与逼近恒等，没有预设一般反演。因此它也提供了以后在测试函数空间上建立变换的起点。

\subsection{\texorpdfstring{\(L^1\)}{L¹} 反演}
\label{b2:subsec:030502}

\begin{theorem}[Fourier 反演]\label{b2:thm:03-inversion}
若 \(f\in L^1(\mathbb R^d)\) 且 \(\widehat f\in L^1(\mathbb R^d)\)，则
\[
 f(x)=\int_{\mathbb R^d}\widehat f(\xi)\cdot e^{2\pi\mathrm{i}x\cdot\xi}\,\mathrm{d}\xi
\]
几乎处处成立。右侧定义一个属于 \(C_0\) 的代表；若 \(f\) 在某一点连续，则该点也成立。
\end{theorem}
\begin{proof}
由控制收敛，\eqref{b2:eq:03-gaussian-inversion}左侧一致趋于 \(g=\check {\widehat f}\)，因为误差不超过
\[
 \int|\widehat f(\xi)|\cdot|e^{-\pi\varepsilon^2|\xi|^2}-1|\,\mathrm{d}\xi\longrightarrow0.
\]
右侧在 \(L^1\) 中趋于 \(f\)。选一列尺度使右侧还有几乎处处收敛子列，便得 \(f=g\) 几乎处处。由 Riemann--Lebesgue 引理，\(g\in C_0\)。若 \(f\) 在 \(x\) 连续，则可取趋向 \(x\)、避开相等关系的零测例外集的一列点；利用两个函数的连续性，得到 \(f(x)=g(x)\)。
\end{proof}

这称为\term[fourier-fanyan-gongshi]{Fourier 反演公式}[Fourier inversion formula]。条件 \(\widehat f\in L^1\) 保证普通反演积分绝对收敛；它不是所有形式的反演都成立的必要条件。尤其对区间示性函数，此条件失败，但仍可讨论带正则化因子的反演。

\subsection{点态反演}
\label{b2:subsec:030503}

沿用第一卷定义 14.4.1：若 \(f\) 在 \(x\) 的小球内满足
\[
 \dfrac1{|B(x,r)|}\int_{B(x,r)}|f(y)-f(x)|\,\mathrm{d}y\longrightarrow0,
\]
则 \(x\) 为该代表的 Lebesgue 点。第一卷命题 14.4.6 证明了非负径向非增单位质量核在 Lebesgue 点处的逼近；本章以 \(f\in L^1\) 满足该命题对非紧支集核要求的全局可积性。Gauss 核和本章 Poisson 核均满足核的条件。

\begin{corollary}[Gauss 正则化反演]\label{b2:cor:03-pointwise-inversion}
若 \(f\in L^1(\mathbb R^d)\)，则在每个 Lebesgue 点 \(x\) 有
\[
 f(x)=\lim_{\varepsilon\downarrow0}
 \int\widehat f(\xi)\cdot e^{-\pi\varepsilon^2|\xi|^2}e^{2\pi\mathrm{i}x\cdot\xi}\,\mathrm{d}\xi.
\]
特别地，该结论几乎处处成立。
\end{corollary}
\begin{proof}
由\eqref{b2:eq:03-gaussian-inversion}，积分就是 \(G_\varepsilon*f(x)\)。应用上述第一卷命题，再用第一卷定理 14.4.2 的几乎处处 Lebesgue 点结论。
\end{proof}

这项结论不要求 \(\widehat f\) 可积，但不能把 Gauss 因子直接删去，也不能未经证明改成任意频率截断形状。一个有两个单侧极限的一维函数，在跳跃点经偶核正则化时趋向中点值：近处分别用左右极限，远处用核在离原点固定距离外的一致衰减估计。点本身的任意赋值不会被积分恢复。

\ProseParagraphBreak

反演还能确定上一节的一维 Poisson 变换。对 \(f(x)=e^{-2\pi t|x|}\)，直接计算已给出 \(\widehat f=P_t\)，且两者均可积。反演与偶性于是给出
\[
 \widehat {P_t}(\xi)=e^{-2\pi t|\xi|}\quad(d=1).
\]
两侧连续，所以最初的几乎处处等式提升为处处等式。

\subsection{唯一性定理}
\label{b2:subsec:030504}

\begin{corollary}[唯一性]\label{b2:cor:03-uniqueness}
若 \(f,g\in L^1(\mathbb R^d)\) 且 \(\widehat f=\widehat g\)，则 \(f=g\) 几乎处处。
\end{corollary}
\begin{proof}
差 \(h=f-g\) 的变换恒为零，特别属于 \(L^1\)，故 Fourier 反演使 \(h=0\) 几乎处处。
\end{proof}

唯一性不依赖一般函数是否存在无正则化的反演积分。它只须对零变换使用反演，所以仍适用于全部 \(L^1\)。例如若 \(f*G_\varepsilon=0\) 对某个固定 \(\varepsilon>0\) 成立，则变换相乘后得到 \(\widehat f(\xi)e^{-\pi\varepsilon^2|\xi|^2}=0\)。Gauss 因子处处非零，因而 \(f=0\)。然而从带误差的卷积数据恢复 \(f\)，却要除以趋零的 Gauss 因子；唯一可恢复与稳定可恢复是不同问题。



\begin{proposition}\label{b2:prop:03-inverse-regularity}
若 \(f\in L^1\)，且 \(\int(1+|\xi|)^m|\widehat f(\xi)|\,\mathrm{d}\xi<\infty\)，其中 \(m\) 为非负整数，则 \(f\) 有一个 \(C^m\) 代表，且其不超过 \(m\) 阶的每个导数都连续并在无穷远趋零。
\end{proposition}
\begin{proof}
零阶权已保证反演适用，取连续代表 \(g=\check {\widehat f}\)。每次对 \(x_j\) 求导，只给频率被积函数增加 \(2\pi\mathrm{i}\xi_j\)。所有相应差商都由假设中的可积矩控制，所以可逐阶求导，得到
\[
 \partial^\alpha g(x)=\int(2\pi\mathrm{i}\xi)^\alpha\widehat f(\xi)e^{2\pi\mathrm{i}x\cdot\xi}\,\mathrm{d}\xi.
\]
其绝对值由频率矩积分控制，连续性与无穷远消失由 Riemann--Lebesgue 引理给出。
\end{proof}

这与空间矩产生频率微分的命题构成对称关系。频率中的可积衰减可以恢复空间光滑性，条件必须作用于带权的绝对积分；仅知道变换逐点趋零，没有提供这些矩。反演因此不仅是一条恢复数值的公式，也是从频率估计提取正则性的工具。


\Cref{b2:tab:03-transform-rules}集中列出本章的运算规则，便于在后续方程计算中同时核对公式与假设。
% 本书自编查阅表；全部公式与条件来自第三章已证明结果。
\begin{table}[H]
\centering\small\renewcommand{\arraystretch}{1.2}\setlength{\tabcolsep}{4pt}
\caption[Fourier 变换的运算规则]{Fourier 变换的运算规则。\(d\) 为正整数，\(x,\xi\in\mathbb R^d\)，统一采用指数 \(e^{-2\pi\mathrm{i}x\cdot\xi}\)；允许复值函数，实值情形包含在内。第三列列出各行的可积性与参数条件。}
\label{b2:tab:03-transform-rules}
\begin{tabularx}{\linewidth}{@{}>{\raggedright\arraybackslash}p{3.4cm}>{\raggedright\arraybackslash}p{6.1cm}>{\raggedright\arraybackslash}X@{}}
\toprule
空间中的操作 & 频率中的表达 & 适用条件 \\
\midrule
\(f(x-h)\) & \(e^{-2\pi\mathrm{i}h\cdot\xi}\cdot\widehat f(\xi)\) & \(f\in L^1(\mathbb R^d)\)，\(h\in\mathbb R^d\) \\
\(e^{2\pi\mathrm{i}\eta\cdot x}\cdot f(x)\) & \(\widehat f(\xi-\eta)\) & \(f\in L^1(\mathbb R^d)\)，\(\eta\in\mathbb R^d\) \\
\(f(Ax)\) & \(|\det A|^{-1}\cdot\widehat f(A^{-\mathsf T}\xi)\) & \(f\in L^1(\mathbb R^d)\)，\(A\in\mathbb R^{d\times d}\)，\(\det A\ne0\) \\
\(\partial_jf\) & \(2\pi\mathrm{i}\xi_j\widehat f(\xi)\) & \(1\leq j\leq d\)，\(f\in L^1(\mathbb R^d)\)，弱导数 \(\partial_jf\in L^1(\mathbb R^d)\) \\
\(x_jf\) & \(\dfrac{\mathrm{i}}{2\pi}\partial_{\xi_j}\widehat f\) & \(1\leq j\leq d\)，\(f,x_jf\in L^1(\mathbb R^d)\) \\
\(f*g\) & \(\widehat f\cdot\widehat g\) & \(f,g\in L^1(\mathbb R^d)\) \\
\(\begin{gathered}G_\varepsilon(x)=\\\varepsilon^{-d}e^{-\pi|x|^2/\varepsilon^2}\end{gathered}\) & \(e^{-\pi\varepsilon^2|\xi|^2}\) & \(\varepsilon>0\)，此核可积且积分为 \(1\) \\
\bottomrule
\end{tabularx}
\end{table}



\begin{chaptersummary}
可积性使 Fourier 变换成为定义在每个频率处的连续函数。平移与调制互换，线性伸缩产生逆转置频率和体积因子；空间矩负责频率微分，空间导数负责频率衰减。卷积在适当指数下存在，其变换是两个变换的乘积。

\ProseParagraphBreak

反演通过逼近恒等取得，而不是交换一个没有收敛依据的双重积分。若变换可积，可以去掉正则化，得到连续代表；若只知道原函数可积，则在 Lebesgue 点保留 Gauss 极限。唯一性仍对整个 \(L^1\) 成立，但高频信息可能在磨光后变得极小，稳定恢复需要额外估计。下一章将以平方能量建立更大的变换空间。
\end{chaptersummary}

\BookExercises

这些题目分别检验相位、卷积的几何形状、逼近条件、误差阶和反演的稳定性。除明确注明的一维题外，定义域均为 \(\mathbb R^d\)。

% 自拟：积分三角不等式的取等条件与调制相联系。
\begin{exercise}\label{b2:ex:03-phase-equality}
设 \(0\ne f\in L^1\)。证明 \(|\widehat f(\xi_0)|=\|f\|_1\) 当且仅当存在 \(|c|=1\)，使 \(f(x)=c e^{2\pi\mathrm{i}x\cdot\xi_0}|f(x)|\) 几乎处处。若 \(f\geq0\)，证明对每个 \(\xi_0\ne0\) 都有严格不等式。
\end{exercise}
\begin{solution}
令 \(h(x)=f(x)\cdot e^{-2\pi\mathrm{i}x\cdot\xi_0}\)。若积分模等于绝对积分，选 \(c\) 使 \(\overline c\int h=\int|h|\)。非负函数 \(|h|-\operatorname{Re}(\overline ch)\) 的积分为零，故它几乎处处为零。这迫使 \(h=c|h|\)，反向由代入得到。若 \(f\geq0\) 且 \(\xi_0\ne0\)，取等要求 \(e^{-2\pi\mathrm{i}x\cdot\xi_0}\) 在 \(\{f>0\}\) 上几乎处处固定。其任一水平集包含于可数个平行超平面，Lebesgue 测度为零；而非零非负可积函数的正集具有正测度，矛盾。
\end{solution}

% 自拟：矩形卷积生成三角函数，以反演计算sinc平方积分。
\begin{exercise}\label{b2:ex:03-tent}
在一维令 \(b=\mathbf1_{[-1/2,1/2]}\)。证明 \(b*b(x)=(1-|x|)_+\)，求其 Fourier 变换，并用反演公式计算
\[
 \int_{\mathbb R}\left(\dfrac{\sin(\pi\xi)}{\pi\xi}\right)^2\,\mathrm{d}\xi.
\]
\end{exercise}
\begin{solution}
卷积等于区间 \([-1/2,1/2]\) 与 \([x-1/2,x+1/2]\) 的重叠长度，所以为三角形函数。区间变换和卷积定理给出其变换为 \((\sin(\pi\xi)/(\pi\xi))^2\)，原点补为 \(1\)。该函数有界，且无穷远以 \(|\xi|^{-2}\) 控制，因而可积。在三角形函数的连续点 \(x=0\) 使用反演，所得积分为 \(1\)。
\end{solution}

% 自拟：逼近恒等的绝对质量界不可由总质量和尾部替代。
\begin{optionalexercise}\label{b2:ex:03-signed-kernel}
在一维设
\[
 K_\varepsilon=G_\varepsilon+\varepsilon^{-1}(\tau_\varepsilon G_\varepsilon-G_\varepsilon),
 \quad0<\varepsilon<1.
\]
证明其积分为 \(1\)，绝对积分尾部在每个固定的原点邻域外趋零，但 \(\|K_\varepsilon\|_1\to\infty\)。再证明，对 \(f\in C_c^2(\mathbb R)\)，有 \(K_\varepsilon*f\to f-f'\) 一致成立。
\end{optionalexercise}
\begin{solution}
平移不改变积分，所以总积分为 \(1\)。作尺度换元，
\[
 \|\tau_\varepsilon G_\varepsilon-G_\varepsilon\|_1
 =\|G(\cdot-1)-G\|_1=c>0,
\]
因此 \(\|K_\varepsilon\|_1\geq c/\varepsilon-1\)。对固定 \(\delta>0\)，当 \(\varepsilon<\delta/2\) 时，两个 Gauss 尾部都由尺度变量中 \(|z|>\delta/(2\varepsilon)\) 的指数尾控制；乘以 \(1+2/\varepsilon\) 后仍趋零。

设 \(h_\varepsilon=G_\varepsilon*f\)，则 \(h_\varepsilon\to f\)、\(h_\varepsilon'\to f'\) 一致成立。又有
\[
 \dfrac{h_\varepsilon(x-\varepsilon)-h_\varepsilon(x)}{\varepsilon}
 =-\int_0^1h_\varepsilon'(x-s\varepsilon)\,\mathrm{d}s\longrightarrow-f'(x)
\]
一致成立，因为 \(f'\) 一致连续。把这一式加到 \(h_\varepsilon\) 上便得到结论。
\end{solution}

% 自拟：利用消失的一阶矩改进逼近误差阶。
\begin{exercise}\label{b2:ex:03-gaussian-error}
设 \(f\in C^2(\mathbb R^d)\) 有界，且 Hesse 矩阵的算子范数至多为 \(M\)。证明
\[
 \|G_\varepsilon*f-f\|_\infty\leq\dfrac{dM}{4\pi}\varepsilon^2.
\]
说明为什么只使用 Lipschitz 条件通常只能得到一次阶误差。
\end{exercise}
\begin{solution}
Taylor 公式给出 \(|f(x-y)-f(x)+\nabla f(x)\cdot y|\leq M|y|^2/2\)。Gauss 核为偶函数，一阶矩为零，因此线性项的积分消失。由一维分部积分，\(\int x^2e^{-\pi x^2}\,\mathrm{d}x=1/(2\pi)\)，所以 \(\int|y|^2G_\varepsilon(y)\,\mathrm{d}y=d\varepsilon^2/(2\pi)\)，得到所求界。只有 Lipschitz 常数 \(L\) 时，直接估计给出 \(L\int|y|G_\varepsilon(y)\,\mathrm{d}y=C_dL\varepsilon\)；缺少可相消的线性近似，不能由这个条件普遍推出二次阶。
\end{solution}

% 自拟：反演与唯一性建立后计算Poisson半群。
\begin{exercise}\label{b2:ex:03-poisson-semigroup}
在一维证明 \(P_t*P_s=P_{t+s}\)，并证明 \(u(t,x)=P_t*f(x)\) 对 \(f\in L^1(\mathbb R)\) 满足 \(u_{tt}+u_{xx}=0\) 于 \(t>0\)，同时 \(u(t,\cdot)\to f\) 于 \(L^1\)。
\end{exercise}
\begin{solution}
三个核都可积。它们的变换分别为 \(e^{-2\pi t|\xi|}\)、\(e^{-2\pi s|\xi|}\)、\(e^{-2\pi(t+s)|\xi|}\)，由卷积定理与唯一性得到半群式。对每个 \(t\geq\varepsilon>0\) 的局部范围，\(P_t\) 的二阶 \(t,x\) 导数都由一个有限常数控制，故可在与 \(f\) 的卷积中求导。直接求导有 \((\partial_t^2+\partial_x^2)P_t(x)=0\)，从而得到方程。初值收敛由逼近恒等定理给出，不要求 \(f\) 有逐点边界值。
\end{solution}

% 自拟：卷积幂等元与频率空间的连通性。
\begin{exercise}\label{b2:ex:03-idempotent}
证明若 \(f\in L^1(\mathbb R^d)\) 满足 \(f*f=f\)，则 \(f=0\) 几乎处处。与圆周上的常函数比较，说明底空间改变后哪一步不同。
\end{exercise}
\begin{solution}
变换给出 \(\widehat f(\xi)^2=\widehat f(\xi)\)，所以每个值为 \(0\) 或 \(1\)。连续函数从连通空间 \(\mathbb R^d\) 映入这两个点，只能为常数；又因无穷远趋零，它恒为零。最后用唯一性。圆周上的频率空间为离散的 \(\mathbb Z\)，系数可以在某些频率上取 \(1\)、其余取 \(0\)；常函数 \(1\) 就满足周期卷积幂等，而不违反任何连通性结论。
\end{solution}

% 自拟：Gauss磨光的单射性与逆操作的不稳定性。
\begin{optionalexercise}\label{b2:ex:03-unstable-deconvolution}
在一维令 \(f_n(x)=e^{-\pi x^2}e^{2\pi\mathrm{i}nx}\)。计算 \(G*f_n\)，证明 \(\|f_n\|_1=1\) 而 \(\|G*f_n\|_1\to0\)。据此说明，即使 Gauss 卷积是单射，也不存在常数 \(C\) 使 \(\|f\|_1\leq C\|G*f\|_1\) 对所有 \(f\in L^1\) 成立。
\end{optionalexercise}
\begin{solution}
在卷积积分中写 \(y=x/2+z\)，平方项满足 \((x-y)^2+y^2=x^2/2+2z^2\)。用 Gauss 的已知变换，得到
\[
 (G*f_n)(x)=2^{-1/2}e^{-\pi n^2/2}e^{-\pi x^2/2}e^{\pi\mathrm{i}nx}.
\]
因此 \(\|G*f_n\|_1=e^{-\pi n^2/2}\)，而 \(\|f_n\|_1=\int G=1\)。若所述逆估计成立，便有 \(1\leq Ce^{-\pi n^2/2}\) 对全部 \(n\) 成立，矛盾。
\end{solution}

% 自拟：卷积支集估计与几何重叠，避免把包含误认为等号。
\begin{exercise}\label{b2:ex:03-support-convolution}
设 \(f,g\in L^1(\mathbb R^d)\) 分别几乎处处支撑在紧集 \(A,B\) 内。证明 \(f*g\) 几乎处处支撑在 \(A+B=\{a+b:a\in A,b\in B\}\) 内。若 \(f,g\) 是非负连续紧支集函数，证明其支集恰为两支集的 Minkowski 和。
\end{exercise}
\begin{solution}
若 \(x\notin A+B\)，则不可能同时有 \(y\in B\) 与 \(x-y\in A\)，所以选取在相应紧集外为零的代表后，卷积积分为零。一般代表给出相同的等价类，故得到第一项。

非负连续情形下，若 \(f(a)>0\)、\(g(b)>0\)，连续性保证在 \(y=b\) 的某个球上同时有 \(g(y)>0\)、\(f(a+b-y)>0\)，因此卷积在 \(a+b\) 严格为正。对支集中的任意 \(a,b\)，用正值点列逼近，得到 \(a+b\) 属于卷积支集。反向已由第一项证明；连续卷积在支集外的开集上若几乎处处为零，便处处为零。非负性在排除抵消时使用，不能省去。
\end{solution}
